\section{Protocol Record and Validation Fields}\label{supp:fields}

This supplement expands the field-level protocol and proofs summarized in
Section III of the main paper. Committed state refers only to the feedback
sub-state $H$ written from generated output.

\subsection{Transaction and certificate records}

For proposal or authority output $\out_t$, the runtime materializes
\[
 \tx_t=(\mathrm{sid},e,p,v_{\mathrm{parent}},[a_s,a_e),\out_t,
 \proj(\out_t;H_t),\mathrm{cert}_t).
\]
The fields carry the session identifier, contract epoch, active policy, parent
version read, covered audio interval, realized output, projected feedback
state, and certificate status. Proposal and authority granularities $g_p$ and
$g_a$ define $\granratio=g_a/g_p\geq1$.

For an authority result $\proj_j^a=\proj(\out_j^a;H_{v_j})$, the protected
coordinator issues
\[
\begin{split}
 \cert_j=\mathrm{Bind}(&\mathrm{sid}_j,e_j,p_j,v_j,[f_j,f_{j+1}),\\
 &h(\audio_{[f_j,f_{j+1})}),h(\proj_j^a),\chi_j).
\end{split}
\]
The configuration digest $\chi_j$ covers the model checkpoint, tokenizer,
decoding parameters, projection-adapter version, numerical precision,
authority-runtime version, and the randomness policy required by the contract.
The audio and projection digests are recomputed over immutable bytes. Each
certificate also carries an unconsumed/consumed state maintained atomically by
the ledger.

\subsection{Atomic validation}

For
\[
 \commit_j=(\mathrm{sid}_j,e_j,p_j,v_j,[f_j,f_{j+1}),
 \proj_j^a,\cert_j),
\]
validation and installation execute as one compare-and-swap transaction. The
runtime checks:
\begin{enumerate}
 \item \emph{session fence}: the session identifier denotes the live ledger;
 \item \emph{contract fence}: epoch and policy equal the active contract;
 \item \emph{parent ownership}: $v_j$ equals the committed state version;
 \item \emph{frontier continuity}: the interval begins at the committed
 frontier and advances it without a gap or overlap;
 \item \emph{content binding}: all certificate fields match the presented
 tuple, and the audio and projection digests recompute;
 \item \emph{single use}: the certificate has not been consumed.
\end{enumerate}
Failure of any field rejects the tuple without changing $H$, the frontier, or
the certificate-consumption bit. Success installs $\proj_j^a$, advances the
frontier and version, and consumes $\cert_j$ atomically.

\section{Detailed Safety Proofs}\label{supp:proofs}

\begin{theorem}[Uncertified-Ancestry Exclusion]\label{supp:thm:ancestry}
If committed state changes exclusively through valid commits, no fine
projection unbound by a valid certificate can enter the ancestry of any future
committed state.
\end{theorem}

\begin{proof}
Index accepted commits by $j$. At initialization, $H_0$ is contract-defined
and contains no generated ancestor. Assume every generated projection on the
unique path to $H_j$ is certificate-bound. The only operation that can create
$H_{j+1}$ is an accepted compare-and-swap. Content validation recomputes
$h(\proj_j^a)$ and requires it to match the projection digest in the presented
certificate; hence different content cannot be installed. A displayed fine
projection has no commit capability and therefore cannot alter the path.

If a fine projection differs from $\proj_j^a$, it fails content binding. If it
happens to be bitwise equal, the installed bytes are nevertheless authorized
by the authority certificate and are attributable to that valid commit, not
to the uncertified proposal. Thus every generated ancestor of $H_{j+1}$ is
certificate-bound, completing the induction.
\end{proof}

\begin{theorem}[Fork-Free Committed-State Linearity]\label{supp:thm:linear}
If validation and installation form one atomic compare-and-swap and each
certificate is consumed at most once, parent-version and frontier checks make
accepted commits fork-free, gap-free, and overlap-free.
\end{theorem}

\begin{proof}
Suppose two distinct candidates read parent version $v$. At most one can win
the compare-and-swap that both validates and installs against $v$. The winner
increments the committed version, so the loser is stale when it reaches the
serialization point and fails the parent check. Therefore no two accepted
children share one committed parent. Single-use consumption additionally
prevents a winning certificate from being replayed at another version.

Let the current committed frontier be $f_j$. A winning interval must begin at
$f_j$ and advances the frontier to $f_{j+1}$. The next accepted interval must
start at this new value. Induction over winners therefore yields the single
sequence $[f_0,f_1),[f_1,f_2),\ldots$, excluding gaps and overlaps. Combining
unique parentage with contiguous intervals yields one linear committed
history.
\end{proof}

\begin{corollary}[Stale-descendant rejection]
A descendant computed from $v_{\mathrm{parent}}\neq\mathrm{version}(H_j)$
cannot commit, even if its output or projection equals current content.
\end{corollary}

\begin{proof}
Parent ownership is a required field check independent of content equality;
the mismatch rejects the compare-and-swap before installation.
\end{proof}

\begin{remark}[Concurrency boundary]
The compare-and-swap is a serialization point for accepted commits. The claim
does not extend to Herlihy--Wing linearizability of all client invocations and
responses, for which a separate real-time concurrent-object specification
would be required.
\end{remark}

\begin{theorem}[Granularity-Independent Safety]\label{supp:thm:granularity}
For every $\granratio\geq1$, Theorems~\ref{supp:thm:ancestry}
and~\ref{supp:thm:linear} hold if each authority segment starts at the
committed parent, covers the next contiguous frontier, produces a legal
projection, and passes certificate and version checks.
\end{theorem}

\begin{proof}
Neither preceding proof counts proposal windows or constrains segment length.
The ancestry proof uses only content binding and the exclusive commit path.
The linear-history proof uses only atomic parent ownership, single-use
consumption, and frontier continuity. Increasing $\granratio$ changes how
many fine regions a certificate covers but changes none of these predicates.
Therefore both invariants are preserved.
\end{proof}

At $\granratio=1$, the proposal and authority regions coincide, directional
coverage is the identity witness, and selective fine replay reduces to
risk-triggered re-verification. This is a protocol reduction only: the
evaluated SG and Coarse/Hybrid contracts use distinct FP16-singleton and
FP32-batched canonical paths and need not be bitwise equal.

\section{Authority-Cost Assumptions}\label{supp:cost}

Let $T_p$ be proposer time, $B$ concurrent sessions, $q$ the SG verification
fraction, $r$ the Hybrid replay fraction, $C_s$ singleton authority cost per
fine region, $C_c(\granratio,B)$ batched coarse cost, $Q_s$ singleton
queueing/fragmentation, and $C_{\mathrm{commit}}$ transaction overhead. The
main-paper decomposition assumes
\[
 T_{\mathrm{service}}\approx\max(T_p,T^a)+T_{\mathrm{repair}}.
\]
Without assuming either lane is critical, Coarse-DGFT improves service time
whenever
\[
\begin{split}
 &\max\!\left(T_p,C_c(\granratio,B)+rB\granratio C_s+
 C_{\mathrm{commit}}\right)+T_{\mathrm{repair}}^D\\
 &\quad<\max\!\left(T_p,qB\granratio C_s+Q_s\right)
 +T_{\mathrm{repair}}^S.
\end{split}
\]
The simpler inequality stated in the main paper is sufficient only when the
authority lane is on or near the critical path in both configurations and
repair costs are comparable. These equations predict operating-point
behavior; they do not identify a precision-independent causal effect because
the evaluated contracts bind different complete authority paths.

\section{Host-Side Protocol Overhead}\label{supp:overhead}

We isolate CPU transaction cost with the production functions in
\texttt{src/revtrace/dgft.py}. Each of 31 randomized-order repeats times
20{,}000 valid calls after 2{,}000 warmup calls per variant. The fixed fixture
uses a 256-token parent prompt, 32 display tokens, and eight speculative
descendants. Table~\ref{supp:tab:overhead} reports percentiles across the 31
repeat-level per-operation aggregates. The host runs Windows~11 and CPython
3.13.5 on a 24-logical-CPU Intel Family~6 Model~183 processor; the review
artifact records the complete environment and source hashes.

\begin{table}[t]
\caption{Host-side production-function microbenchmark ($\mu$s/operation).}
\label{supp:tab:overhead}
\centering
\footnotesize
\begin{tabular}{@{}lrrr@{}}
\toprule
Path & Median & P95 & P99 \\
\midrule
State materialization; no validation & 19.077 & 19.813 & 19.885 \\
Job/owner/version/frontier validation & 24.835 & 25.636 & 25.887 \\
Job plus full certificate validation & 80.604 & 82.902 & 85.866 \\
Production commit & 101.107 & 104.486 & 104.514 \\
\bottomrule
\end{tabular}
\end{table}

The paired median increment from no-check materialization to the complete
production commit is 81.960\,$\mu$s. In 1{,}048 existing online commits, the
broader coordinator critical section has median/P95/P99
1.318/2.717/3.784\,ms; that interval also includes lock wait, conversion,
certificate construction and hashing, materialization, and pending-version
updates, so it is only a deployed-path upper envelope. Neither measurement is
used as GPU or end-to-end speedup evidence.
